% ============================================================
%  abstract.tex — English Abstract
% ============================================================
\chapter*{Abstract}
\addcontentsline{toc}{chapter}{Abstract}

Breast cancer molecular subtype classification from
Dynamic Contrast-Enhanced MRI (DCE-MRI) is a clinically
important task: the subtype determines whether a patient
is eligible for hormone therapy or requires chemotherapy.
Training an automated classifier is, however, difficult
in practice.
A single hospital typically holds fewer than 800 patients,
the class imbalance between subtypes reaches 10:1, and
data protection regulations --- including Algerian Law
18-07 --- prevent institutions from pooling their imaging
archives.

This work investigates whether federated learning can
produce a useful classifier under these constraints,
without any patient data leaving its source institution.
The study is conducted on a public 737-patient breast
DCE-MRI cohort (Duke Hospital, Zenodo DOI
10.5281/zenodo.7956360) with four molecular subtypes:
Luminal~A, Luminal~B, HER2-enriched, and Triple-Negative.

A three-stage local training pipeline was built from
validated components: a ConvNeXt-Nano backbone with
Gated Attention Multiple Instance Learning processes
whole volumes without requiring tumour segmentation masks;
Label-Distribution-Aware Margin loss, Adaptive
Sharpness-Aware Minimisation, Stochastic Weight Averaging,
and Classifier Retraining address class imbalance at
each stage of training.
This pipeline achieves a four-class ensemble macro~F1
of~0.342 and a binary macro~F1 of~0.667 on the
held-out validation fold.

For the federated experiment, three simulated hospital
clients are constructed from a Dirichlet partition
($\alpha = 0.5$), producing deliberately heterogeneous
local class distributions.
Standard FedAvg collapses to majority-class prediction
(macro~F1~= 0.429).
SCAFFOLD improves early representation learning but fails
at head calibration under the extreme class asymmetry.
The proposed FedSCRT algorithm --- which combines a
SCAFFOLD backbone with federated class-balanced classifier
retraining --- achieves macro~F1~= 0.629 and
AUC-ROC~= 0.687, closing 84.1\% of the federated
performance gap and exceeding the centralised AUC baseline.
A theoretical result establishes that class-balanced
sampling renders each client's head gradient independent
of its local class prior, and a controlled cosine-similarity
experiment on three artificial hospitals confirms this:
balanced gradients align with the global objective
($\cos \geq 0.99$) regardless of local prior, while naturally
sampled gradients drift toward the majority class, explaining
why the federated classifier-retraining stage is robust to the
hospitals' heterogeneous class distributions.

\vspace{1em}
\noindent\textbf{Keywords:} Federated learning, breast MRI,
molecular subtype classification, class imbalance, FedSCRT,
Gated Attention MIL, LDAM, SCAFFOLD, USTHB.
